%\subsection{Proof of \cref{lem:completelattice}}
% \begin{lemma}[Restatement of Lemma~\ref{lem:completelattice}]
% 	$\Ucal_n$ is a complete lattice and,
% 	for any $\phi, \xi$,
% 	the function $\lambda f.\lambda \vec b.\sem[{\xi[X\mapsto f]}]{\phi\{\vec b / \vec x\}}$ is monotone.
% \end{lemma}
\nomlemma*

\begin{proof}
	\begin{enumerate}
	\item For the former, we have that 
	$$\orb(\cls{}(X)) = \{\orbs(\pi\cdot x) \mid x \in X, \pi \in \Perm\} = \{\orbs(x) \mid x \in X\}$$
	For the latter, we define the function $f: \orb(\cls{}(X)) \to \orb[\supp(X)](X)$ as follows:
	\[
	f(O) = \orbs[\supp(X)](y) \text{ for some } y \in O \cap X
	\]
	Let us choose $O_1, O_2 \in \orb(\cls{}(X))$ such that $O_1 \neq O_2$ and $\hat x_1 \in O_1 \cap X, \hat x_2 \in O_2 \cap X$.
	Suppose $f(O_1) = \orbs[\supp(X)](\hat x_1) = \orbs[\supp(X)](\hat x_2) = f(O_2)$. 
	This means that $\hat x_1, \hat x_2 \in \orbs[\supp(X)](\hat x_1)$,
	and so there is some $\pi \in \fix(\supp(X))$ such that $\pi \cdot \hat x_1 = \hat x_2$. 
	As $\fix(\supp(X)) \subseteq \Perm$, this implies $\hat x_2 \in O_1$, which is a contradiction. 
	Thus, $f$ is a 1-1 function and $|\orb(\cls{}(X))| \le |\orb[\supp(X)](X)|$ as required.	
	
	\item We have that 
	\begin{align*}
		\cls{}(X \times Y) &= \{\pi \cdot (x, y) \mid \pi \in \Perm, (x,y) \in X \times Y\}
		\\ &= \{(\pi \cdot x, \pi \cdot y) \mid \pi \in \Perm, (x, y) \in X \times Y\}
	\end{align*} 
	As $\supp(X) = \emptyset$, this means that $X$ is closed under permutation, and so we obtain:
	\begin{align*}
		\{(\pi \cdot x, \pi \cdot y) \mid \pi \in \Perm, (x, y) \in X \times Y\} &= \{(x, \pi \cdot y) \mid \pi \in \Perm, (x, y) \in X \times Y\}\\
		&= X \times \{\pi \cdot y \mid \pi \in \Perm, y \in Y\} \\ &= X\times \clos(Y)
	\end{align*}
%	3
	\item Let $x_{O}$ range over orbit representatives in $\orb(\Xcal), \orb(\Ycal)$ such that for any $O_\Xcal \in \orb(\Xcal)$, $x_{O_\Xcal} \in O_\Xcal$, and similarly for $\Ycal$.
	We define the following function:
	\[
	F(O_\Xcal, O_\Ycal) = \bigcup_{\pi' \in \Perm}\{\pi \cdot (x_{O_\Xcal}, \pi'\cdot x_{O_\Ycal}) \mid \pi \in \Perm\}
	\]
	Where $O_\Xcal \in \orb(\Xcal)$ and $O_\Ycal \in \orb(\Ycal)$. 
	Without loss of generality, let us assume that $\supp(x_{O_\Xcal}) \cap \supp(x_{O_\Ycal}) = \emptyset$. 
	Then, we can say that $\{\pi \circ \pi' \mid \pi, \pi' \in \Perm\} = \{\pi \circ \pi' \mid \pi \in \Perm, \pi' \in \fix(\Atoms \setminus \supp(x_{O_\Xcal}, x_{O_\Ycal})\}$ as both sets are closed under permutation.  
	Let us fix $\vec b_0$ to be some sequencing of $\supp(x_{O_\Ycal})$ and let us write $\vec b' \lesssim \vec b_0$ to indicate $\vec b'$ is a sub-sequence of $\vec b_0$.
	From this, we can say:
	\begin{align*}
		F(O_\Xcal, O_\Ycal) &= \bigcup_{\pi' \in \fix(\Atoms \setminus \supp(x_{O_\Xcal}, x_{O_\Ycal}))}\{\pi \cdot (x_{O_\Xcal}, \pi'\cdot x_{O_\Ycal}) \mid \pi \in \Perm\} \\
		&= \bigcup_{\substack{\vec b_0 \lesssim \vec b,\\ \{\vec a\} \subseteq \supp(x_{O_\Xcal})}} \{\pi \cdot (x_{O_\Xcal}, \move{\vec b}{\vec a}\cdot x_{O_\Ycal}) \mid \pi \in \Perm\} \\
		&= \bigcup_{\substack{\vec b_0 \lesssim \vec b,\\ \{\vec a\} \subseteq \supp(x_{O_\Xcal})}} \orbs((x_{O_\Xcal}, \move{\vec b}{\vec a}\cdot x_{O_\Ycal}))
	\end{align*}
	Next we can write $\orb(\Xcal \times \Ycal)$ in terms of $F$ as follows:
	\begin{align*}
		\orb(\Xcal\times \Ycal) &= \{\{(\pi \cdot x_{O_\Xcal},\pi'\cdot x_{O_\Ycal}) \mid \pi, \pi' \in \Perm\} \mid O_\Xcal \in \orb(\Xcal), O_\Ycal \in \orb(\Ycal)\}\\
		&= \{\{\pi\cdot (x_{O_\Xcal},\pi^{-1}\cdot \pi'\cdot x_{O_\Ycal}) \mid \pi, \pi' \in \Perm\} \mid O_\Xcal \in \orb(\Xcal), O_\Ycal \in \orb(\Ycal)\}\\
		&= \{\{\pi\cdot (x_{O_\Xcal},\pi''\cdot x_{O_\Ycal}) \mid \pi, \pi'' \in \Perm\} \mid O_\Xcal \in \orb(\Xcal), O_\Ycal \in \orb(\Ycal)\}\\
		&= \bigcup_{\substack{O_\Xcal \in \orb(\Xcal),\\ O_\Ycal \in \orb(\Ycal)}}F(O_\Xcal, O_\Ycal)\\
		\therefore\ |\orb(\Xcal \times \Ycal)| &\le \sum_{\substack{O_\Xcal \in \orb(\Xcal),\\ O_\Ycal \in \orb(\Ycal)}}|F(O_\Xcal, O_\Ycal)| 
	\end{align*}
	It remains to calculate $|F(O_\Xcal, O_\Ycal)|$ for any $O_\Xcal \in \orb(\Xcal), O_\Ycal \in \orb(\Ycal)$. 
	Let us write $m_1 = |\supp(x_{O_\Xcal})|$ and $m_2 = |\supp(x_{O_\Ycal})|$,  assuming WLOG that $m_1 \ge m_2$.
	% for any $X_A \in O_A,X_B \in O_B$ and $\pi' \in \fix(\Atoms \setminus \supp(X_A, X_B))$ we have:
	\begin{align*}
		|F(O_\Xcal, O_\Ycal)| &= |\bigcup_{\substack{\vec b \lesssim \vec b_0,\\ \{\vec a\} \subseteq \supp(x_{O_\Xcal})}}\orbs((x_{O_\Xcal}, \move{\vec b}{\vec a}\cdot x_{O_\Ycal}))|\\
		&=
		|\{\move{\vec b}{\vec a} \mid \{\vec a\} \subseteq \supp(x_{O_\Xcal}), \vec b \lesssim \vec b_0\}|
		\\&= |\{\sigma: \supp(x_{O_\Ycal}) \rightharpoonup \supp(x_{O_\Xcal}) \mid \sigma\text{ is a 1-1 function}\}|
		\\ &= m_1(m_1 - 1)\dots(m_1 - m_2 + 1) 
		\\&+ (m_2)(m_1)(m_1 - 1)\dots(m_1 - m_2 + 2) + \dots + 1
		\\ &=\sum^{m_2}_{i = 0} \choice{m_2}{i} m_1(m_1 - 1)\dots(m_1 - i + 1)  
		% \\&\le \sum^{n_\Ycal}_{i = 0} 
		% \choice{n_\Ycal}{i}
		% (n_\Xcal-1)^i \cdot \frac{n_\Xcal}{n_\Xcal - 1}
		% \\ &\le {n_\Xcal}^{n_\Ycal} \cdot \frac{n_\Xcal}{n_\Xcal-1} \le {n_\Xcal}^{n_\Ycal}(1 + \epsilon) 
	\end{align*}
        Let us call the above expression $A(m_1,m_2)$.
        Suppose $\{n_1, n_2\} = \{\regindex(\Xcal), \regindex(\Ycal)\}$ such that $n_1 \ge n_2$.
        Then we have:
        \[
        |\orb(\Xcal \times \Ycal)| \le |\orb(\Xcal)|\cdot|\orb(\Ycal)|\cdot A(n_1,n_2)
      \]
      Moreover:
      \begin{align*}
A(n_1,n_2) &\le \sum^{n_2}_{i = 0} 
		\choice{n_2}{i}
		(n_1-1)^i \cdot \frac{n_1}{n_1 - 1}
        \\ &
=  \frac{n_1}{n_1 - 1}\cdot\sum^{n_2}_{i = 0} 
		\choice{n_2}{i}
		(n_1-1)^i 
=  \frac{n_1}{n_1 - 1}\cdot {n_1}^{n_2}=  {n_1}^{n_2}\cdot (1 + \epsilon) 
      \end{align*}
      with $\epsilon=1/(n_1-1)$. We finally show that a bound without $\epsilon$ can be obtained when $n_2\ge 4$. Let us assume that $n_2\ge 4$ and set
      $B(n_1,n_2)=\sum^{n_2}_{i = 0} \choice{n_2}{i} (n_1-1)^i$. We break the sums of $A$ and $B$ as follows:
      \begin{align*}
        A(n_1,n_2) = A'(n_1,n_2)+A''(n_1,n_2) &=
\sum^{4}_{i = 0} \choice{n_2}{i} n_1(n_1 - 1)\dots(n_1 - i + 1)\\ & +
\sum^{n_2}_{i = 5} \choice{n_2}{i} n_1(n_1 - 1)\dots(n_1 - i + 1)\\
B(n_1,n_2) = B'(n_1,n_2)+B''(n_1,n_2) &=
\sum^{4}_{i = 0} \choice{n_2}{i} (n_1 - 1)^i+
\sum^{n_2}_{i = 5} \choice{n_2}{i} (n_1 - 1)^i
\end{align*}                                                
Note first that $A''(n_1,n_2)\leq B''(n_1,n_2)$ as:
\begin{align*}
& n_1(n_1-1)(n_1-2)(n_1-3)(n_1-4)\leq (n_1-1)^5\\
&\iff n_1(n_1-2)(n_1-3)(n_1-4)\leq (n_1-1)^4 \\
&\iff n_1^4-9n_1^3+26n_1^2-24n_1\leq n_1^4-4n_1^3+6n_1^2-4n_1+1\\
&\iff -5n_1^3+20n_1^2-20n_1-1\leq 0
\end{align*}
Next, writing $\kappa_i=\choice{n_2}{i}$, we have:
\begin{align*}
  A'(n_1,n_2) &= 1+n_2n_1+\kappa_2n_1(n_1-1)+\kappa_3n_1(n_1-1)(n_1-2)+\kappa_4n_1(n_1-1)(n_1-2)(n_1-3)\\
 B'(n_1,n_2) &= 1+n_2(n_1-1)+\kappa_2(n_1-1)^2+\kappa_3(n_1-1)^3+\kappa_4(n_1-1)^4
\end{align*}
and, hence:
\begin{align*}
  A'(n_1,n_2)-B'(n_1,n_2) &= n_2+\kappa_2(n_1-1)+\kappa_3(-n_1+1)+\kappa_4(-2n_1^3+5n_1^2-2n_1-1)\\
  &= -2\kappa_4n_1^3+5\kappa_4n_1^2+(\kappa_2-\kappa_3-2\kappa_4)n_1+(n_2-\kappa_2+\kappa_3-\kappa_4)
\end{align*}
so, writing $\kappa'_i=\choice{n_2}{i}/n_2$, we have $A'(n_1,n_2)-B'(n_1,n_2)\leq 0$ iff
\begin{align*}
D(n_1,n_2) = -2\kappa_4'n_1^3+5\kappa_4'n_1^2+(\kappa_2'-\kappa_3'-2\kappa_4')n_1+(1-\kappa_2'+\kappa_3'-\kappa_4')\leq0
\end{align*}
Differentiating $D(n_1,n_2)$ for $n_1$ we get:
\begin{align*}
\frac{\partial D}{\partial n_1}= -6\kappa_4'n_1^2+10\kappa_4'n_1+(\kappa_2'-\kappa_3'-2\kappa_4')
\end{align*}
By some elementary analysis we can see that the above is less or equal to 0 when $n_1,n_2\geq 4$. Given that $n_1\geq n_2$, it suffices to show that $D(n_2,n_2)\leq 0$ for $n_2\geq 4$. For the latter, we note that $D(4,4)< 0$ and that $D(n_2,n_2)$ is decreasing for $n\leq4$.
\end{enumerate}
\end{proof}


\completelattice*
\begin{proof}
	First, we show that $\Ucal_n$ is a complete lattice. 
	It suffices to show that for any $S\subseteq \Ucal_n$, $\bigsqcup S$ is the least lower bound of $S$, and $\bigsqcap S$ is the greatest upper bound of $S$. In order to show that $\bigsqcup S$ is the least upper bound of $S$, we must show that:
 \begin{enumerate}
		\item $\forall f\in S$. $f\sqsubseteq \bigsqcup S$
		\item $\forall g\in \Ucal_n.\ (\forall f\in S. f\sqsubseteq g)\implies\bigsqcup S \sqsubseteq g$
	\end{enumerate}
 	For case 1 with $f\in S$, we have that $\forall \vec a\in\Atoms. f(\vec a) \subseteq (\bigsqcup S)(\vec a)$
	and hence $f\sqsubseteq \bigsqcup S$.\\
	%
  	For case 2, let us pick such a $g$. It follows that for all $f\in S$ and $\vec a$ we have that $f(\vec a) \subseteq g(\vec a)$,
	meaning we have $g(\vec a) \supseteq \bigcup \{f(\vec a) \mid f \in S\}$.
	Therefore, $\bigsqcup S \sqsubseteq g$.
	To show that $\bigsqcap S$ is the greatest lower bound, we must show that:
	\begin{enumerate}
		\item $\forall f\in S$. $f\sqsupseteq \bigsqcap S$
		\item $\forall g\in \Ucal_n.\ (\forall f\in S. f\sqsupseteq g)\implies\bigsqcap S \sqsupseteq g$
	\end{enumerate}
	For case 1 with $f\in S$, we have that $\forall \vec a\in\Atoms. f(\vec a)\supseteq (\bigsqcap S)(\vec a)$ and hence $f\sqsupseteq \bigsqcap S$.\\
	For case 2, let us pick such a $g$. It follows that for all $f \in S$ and $\vec a$ we have that $f(\vec a)\supseteq g(\vec a)$, meaning we have $g(\vec a) \subseteq \bigcap\{(f(\vec a)\mid f \in S\}$. Therefore, $\bigsqcap S \sqsupseteq g$.     
	With a least upper bound (the join) and a greatest lower bound (the meet), we conclude that $\Ucal_n$ is a complete lattice. 
	
	Next, we show that the function $\lambda f.\lambda \vec x.\sem[{\xi[X\mapsto f]}]{\phi}$ is monotone.
	For any $\phi',\xi', \vec y, Y$, we define:
	\[\F[\xi']{\phi'}{\vec y}{Y} = \lambda f.\lambda \vec y.\sem[{\xi'[Y\mapsto f]}]{\phi}\]
	 We need to show that for all $\vec c \in \Atoms^{|\vec x|}$
	and $f, f'\in \Atoms^{|\vec x|} \rightarrow \Ucal$ such that $f \sqsubseteq f'$: 
	\begin{align*}
		\F{\phi}{\vec x}{X}(f)(\vec c) \subseteq \F{\phi}{\vec x}{X}(f')(\vec c) &\text{ when $X$ has an even number of negations}\\
		\F{\phi}{\vec x}{X}(f)(\vec c) \supseteq \F{\phi}{\vec x}{X}(f')(\vec c) &\text{ when $X$ has an odd number of negations}
	\end{align*}
	We perform an induction on $|\phi|$
	and state our inductive hypotheses (denoted \textit{IH}) for every $\phi'$ with $|\phi'| < |\phi|$ as follows:
	\begin{enumerate}
		\item $\F{\phi'}{\vec x}{X}(f) \sqsubseteq \F{\phi'}{\vec x}{X}(f')$ ($\phi$ has an even
		number of negations)
		\item $\F{\phi'}{\vec x}{X}(f) \sqsupseteq \F{\phi'}{\vec x}{X}(f')$ ($\phi$ has a odd number of negations)
	\end{enumerate}
	First, we examine the case where $X$ is not present in $\phi$. For any $\vec c \in \Atoms^{|\vec x|}$:
	\[
	\F{\phi}{\vec x}{X}(f)(\vec c) = \F{\phi}{\vec x}{X}(f')(\vec c) = \sem{\phi\{\vec c / \vec x\}}
	\]
	We now examine all cases where $X$ is present in $\phi$, with focus on the cases where every occurrence of $X$ is within an even number of negations. 
	The base case is when $\phi = X(\vec a)$:
	\begin{align*}
		\F{\phi}{\vec x}{X}(f)(\vec a) &= \xi(X)(\vec a) = f(\vec a)\\
		\F{\phi}{\vec x}{X}(f')(\vec a) &= \xi(X)(\vec a) = f'(\vec a)
	\end{align*}
	As $f\sqsubseteq f'$, we have that $f(\vec a) \subseteq f'(\vec a)$ for all $\vec a \in \Atoms^{|\vec x|}$, 
	meaning $\F{\phi}{\vec x}{X}(f) \sqsubseteq \F{\phi}{\vec x}{X}(f')$.\\
	Next we examine inductive steps for all $\vec c \in \Atoms^{|\vec x|}$:
	\begin{itemize}
		\item $\phi = \lnot \phi'$: for $\hat f \in \{f, f'\}$, we have
		\begin{align*}
			&\F{\phi}{\vec x}{X}(\hat  f)(\vec c)  = \fstates \setminus \sem[{\xi[X \mapsto \hat f]}]{\phi'(\vec c)} =  \fstates \setminus \F{\phi'}{\vec x}{X}(\hat f)(\vec c)
		\end{align*}
		By the inductive hypothesis, we have that $\F{\phi'}{\vec x}{X}(f) \sqsupseteq \F{\phi'}{\vec x}{X}(f')$, meaning for all $\vec c \in \Atoms^{|\vec x|}$, we have that
		\begin{align*}
			\gap\F{\phi'}{\vec x}{X}(f)(\vec c) \supseteq \F{\phi'}{\vec x}{X}(f')(\vec c) \\
			\therefore\gap\F{\phi}{\vec x}{X}(f)(\vec c) \subseteq \F{\phi}{\vec x}{X}(f')(\vec c)
		\end{align*}
		as required.
		Note we use the second inductive hypothesis as $\lnot \phi'$ has one more negation than $\phi'$.
		%
		\item $\phi = \phi_1 \lor \phi_2$: for $\hat f \in \{f, f'\}$, we have
		\[
		\F{\phi}{\vec x}{X}(\hat f)(\vec c) = \sem[{\xi[X\mapsto \hat f]}]{\phi_1\{\vec c / \vec x\})} \cup \sem[{\xi[X\mapsto \hat f]}]{\phi_2\{\vec c / \vec x\}} = \F{\phi_1}{\vec x}{X}(\hat f)(\vec c) \cup \F{\phi_2}{\vec x}{X}(\hat f)(\vec c)
		\]
		By the inductive hypothesis, we have that $\F{\phi_i}{\vec x}{X}(f)(\vec c) \subseteq \F{\phi_i}{\vec x}{X}(f')(\vec c)$ for $i \in \{1,2\}$,
		therefore
		$\F{\phi}{\vec x}{X}(f)(\vec c)\subseteq \F{\phi}{\vec x}{X}(f')(\vec c)$.
		\item $\phi = \bigvee_{y} \phi'$: for $\hat f \in \{f,f'\}$, we have
		\[
		\F{\phi}{\vec x}{X}(\hat f)(\vec c) = \bigcup_{a\in \Atoms}\sem[{\xi[X\mapsto \hat f]}]{\phi'\{a/y\}\{\vec c / \vec x\}} = \bigcup_{a \in \Atoms}\F{\phi'\{a/y\}}{\vec x}{X}(\hat f)(\vec c)
		\]
		As for all $a$, $|\phi| > |\phi'\{a/y\}|$, we can use the inductive hypothesis, therefore:    \[
		\bigcup_{a \in \Atoms}\F{\phi'\{a/y\}}{\vec x}{X}(f)(\vec c) \subseteq  \bigcup_{a \in \Atoms}\F{\phi'\{a/y\}}{\vec x}{X}(f')(\vec c)
		\]
		as required.
	    \item $\phi = \fresh{y} \phi'$: for $\hat f \in \{f,f'\}$, we have
		\begin{align*}
			\F{\phi}{\vec x}{X}(\hat f)(\vec c) &= \bigcup_{a\in \Atoms}\{(s, H) \in \sem[{\xi[X\mapsto \hat f]}]{\phi'\{a/y\}\{\vec c/\vec x\}} \mid a \notin H\} 
			\\ &= \bigcup_{a \in \Atoms}\{(s, H) \in \F{\phi'\{a/y\}}{\vec x}{X}(\hat f)(\vec c) \mid a \notin H\}
		\end{align*}
		As for all $a$, $|\phi| > |\phi'\{a/y\}|$, we can use the inductive hypothesis.
		Using the inductive hypothesis, if $(s, H) \in \F{\phi'\{a/y\}}{\vec x}{X}(f)(\vec c)$ then $(s, H) \in \F{\phi'\{a/y\}}{\vec x}{X}(f')(\vec c)$, thus we can conclude.
		\item $\phi = \diam{\ell} \phi'$: for $\hat f \in \{f,f'\}$, we have:
		\begin{align*}
		&\F{\phi}{\vec x}{X}(\hat f)(\vec c) = \{(s, H) \in \fstates \mid \exists s\xrightarrow{\ell}(s', H').(s', H')\in \sem[{\xi[X\mapsto \hat f]}]{\phi'\{\vec c / \vec x\}} \} 
		\\=\ & \{(s', H') \in \fstates \mid \exists s\xrightarrow{\ell}(s', H').(s', H') \in \F{\phi'}{\vec x}{X}(\hat f)(\vec c)\}
		\end{align*}
		Using the inductive hypothesis, if $(s', H')\in \F{\phi'}{\vec x}{X}(f)(\vec c)$ then $(s', H')\in \F{\phi'}{\vec x}{X}(f')(\vec c)$,
		thus, we can conclude.
		\item $\phi = (\mu Y(\vec y).\phi')(\vec a)$: for $\hat f \in \{f,f'\}$, we have:
    	\begin{align*}
			\F{\phi}{\vec x}{X}(\hat f)(\vec c) =\gap (\lfp(\lambda g. \lambda \vec d.\sem[{\xi[X\mapsto \hat f][Y\mapsto g]}]{\phi'(\vec c)\{\vec d/\vec y\}}))(\vec a) \\
			% \lfp(\F[{\xi[X\mapsto \hat f]}]{\phi'(\vec c)}{\vec y}{Y})=
			=\gap \bigsqcap \{g\in \Ucal_n\mid 
			\lambda \vec d. \sem[{\xi[X\mapsto \hat f][Y\mapsto g]}]{\phi'(\vec c)\{\vec d/\vec y\}}\sqsubseteq g\}(\vec a)\\
			=\gap\bigsqcap \{g\in \Ucal_n\mid 
			\lambda \vec d. \sem[{\xi[X\mapsto \hat f][Y\mapsto g]}]{\phi'\{\vec d/\vec y\}(\vec c)}\sqsubseteq g\}(\vec a)\\
			=\gap \bigsqcap\{g \in \Ucal_n\mid \lambda \vec d.\F[{\xi[Y\mapsto g]}]{\phi'\{\vec d/\vec y\}}{\vec x}{X}(\hat f)(\vec c) \sqsubseteq g\}(\vec a)\\
			=\gap \bigsqcap G_{\hat f, \vec c}(\vec a) &[\text{for economy}]
		\end{align*}
		Let us take some $ g' \in G_{f', \vec c}$.
		We have that $\lambda \vec y.\F[{\xi[Y\mapsto g']}]{\phi'}{\vec x}{X}(f')(\vec c) \sqsubseteq g'$ by definition.
		Furthermore, by the inductive hypothesis,
		we have that 
		$$\lambda \vec y.\F[{\xi[Y\mapsto g']}]{\phi'}{\vec x}{X}(f)(c)\sqsubseteq\lambda \vec y.\F[{\xi[Y\mapsto g']}]{\phi'}{\vec x}{X}(f')(c) \sqsubseteq g'$$
		Hence, $g' \in G_{f, \vec c}$ and so: 
		\[
			G_{f', \vec c} \subseteq G_{f, \vec c} \text{ and therefore } \bigsqcap G_{f, \vec c} (\vec a) \subseteq \bigsqcap G_{f', \vec c} (\vec a)
		\]
		as required.
	\end{itemize}
%	\begin{itemize}

		
%	\end{itemize}
\end{proof}
	
\subsection{Proof of \cref{lem:semanticsequivariant}}
\semanticsequivariant*
%	The function $\lambda \phi, \xi.\sem{\phi}$ is equivariant.
%	Therefore, 
%	for any formula $\phi$ and $\Ucal$-variable assignment $\xi$, $\supp(\sem{\phi}) \subseteq \supp(\phi) \cup \supp(\xi)$.
%\end{lemma}
\begin{proof}
	We perform proof by induction on $\phi$.
	Let $\pi$ be some permutation, we show that
	$\semp{\phi} = \pi\cdot(\sem{\phi})$.
	The bases cases are as follows:
	\begin{itemize}
		\item $a=b: \semp{(a=b)} = \sem[\pi\cdot \xi]{\pi(a) = \pi(b)} = \emptyset = \pi\cdot \emptyset = \pi\cdot\sem{a=b}$.
		\item $a=a: \semp{(a=a)} = \sem[\pi\cdot \xi]{\pi(a) = \pi(a)} = \fstates = \pi\cdot \fstates = \pi\cdot\sem{a=a}$.
		\item $X(\vec a): \semp{(X(\vec a))} = \sem[\pi\cdot\xi]{X(\pi\cdot \vec a)} = (\pi\cdot\xi)(X)(\pi\cdot \vec a)$.
		We have that $\pi\cdot\xi = \{\pi\cdot (Y, U)\mid (Y,U)\in \xi\} = \{(Y, \pi \cdot U)\mid (Y,U)\in \xi\}$. 
		Let us say $\xi(X) = U_x$, then $(\pi\cdot\xi)(X) = \pi\cdot U_x$. 
		From this we can say that $\pi\cdot U_x = \{\pi\cdot (\vec b, S)\mid (\vec b, S) \in U_x\}=\{(\pi\cdot \vec b, \pi\cdot S)\mid (\vec b, S) \in U_x\}$.
		It follows that $(\pi\cdot U_x)(\pi\cdot \vec a) = \pi\cdot S$
		which is equal to 
		$\pi\cdot (\xi(X)(\vec a)) = \pi\cdot \sem{X(\vec a)}$.
	\end{itemize}
	We now examine inductive steps.\
	\begin{itemize}
		\item $\phi_1 \vee \phi_2$:
		we have that $\sem[\pi\cdot \xi]{\pi\cdot \phi_1 \vee \pi\cdot \phi_2} = \semp{\phi_1} \cup \semp{\phi_2}$, by the inductive hypothesis, the latter is $\pi\cdot \sem{\phi_1} \cup \pi\cdot \sem{\phi_2}=\pi\cdot (\sem{\phi_1}\cup\sem{\phi_2})$ as required. 
		\item $\lnot \phi_1$: 
		By definition, $\semp{(\lnot\phi_1)} = \fstates\setminus\semp{\phi_1}$, which by applying the inductive hypothesis gives $\fstates\setminus\pi\cdot (\sem{\phi_1})$.
		As $\fstates = \pi\cdot \fstates$, this is equal to $\pi\cdot \fstates\setminus\pi\cdot (\sem{\phi_1})=\pi\cdot(\fstates\setminus(\sem{\phi_1}))=\pi\cdot(\sem{\lnot\phi_1})$. 
		\item $\phi=\bigvee_{x}\phi_1$:
		In this case, $\semp{(\bigvee_{x}\phi_1)} = \bigcup_{a\in\Atoms}\semp{(\phi_1\{a/x\})}$. By inductive hypothesis, the latter is $\bigcup_{a\in\Atoms}\pi\cdot (\sem{(\phi_1\{a/x\})})=\pi\cdot (\bigcup_{a\in\Atoms}\sem{(\phi_1\{a/x\}})=\pi\cdot (\sem{\bigvee_{x}\phi_1})$.
		\item $\phi=\fresh{x}\phi'$:
		In this case, $\semp{(\fresh{x}\phi')} = \bigcup_{a\in\Atoms}\{(s, H) \in \semp{(\phi'\{a/x\})} \mid a \notin H\}$. 
		By inductive hypothesis, the latter is $\bigcup_{a\in\Atoms}\{(s, H) \in \pi \cdot \sem{(\phi'\{a/x\})} \mid a \notin H\}
		% =\bigcup_{a\in\Atoms}\{\pi \cdot (s, H) \in \sem{(\phi'\{a/x\})} \mid \pi(a) \notin H\}
		=\pi \cdot \bigcup_{a\in\Atoms}\{(s, H) \in \sem{(\phi'\{a/x\})} \mid a \notin H\}$.
		\item $\diam{\ell}\phi_1:$ we have that
		\begin{align*}
			\semp{(\diam{\ell}\phi_1)} &= \{s\in\fstates\mid\exists s\xrightarrow{\pi\cdot\ell} s' . s' \in \semp{\phi_1}\}
			= \{s\in\fstates\mid\exists s\xrightarrow{\pi\cdot\ell} s' . s' \in \pi\cdot(\sem{\phi_1})\}\\
			&= \{s\in\fstates\mid\exists s\xrightarrow{\pi\cdot\ell} \pi\cdot s' . \pi\cdot s' \in \pi\cdot(\sem{\phi_1})\}\ (\text{Rewrite}\ [s' = \pi\cdot s'])\\ 
			&= \{s\in\fstates\mid\exists s\xrightarrow{\pi\cdot\ell} \pi\cdot s' . s' \in \sem{\phi_1}\}
			= \{s\in\fstates\mid\exists \pi^{-1}\cdot s\xrightarrow{\ell} s' . s' \in \sem{\phi_1}\}\\
			&= \{\pi\cdot s\in\fstates\mid\exists s\xrightarrow{\ell} s' . s' \in \sem{\phi_1}\}
			= \pi\cdot\{s\in\fstates\mid\exists s\xrightarrow{\ell} s' . s' \in \sem{\phi_1}\}\\
			&= \pi\cdot\sem{(\diam{\ell}\phi_1)}
		\end{align*}
		\item $(\mu X(\vec x).\phi)(\vec a):$ we have that
    \begin{align*}
	\pi\cdot \sem{(\mu X(\vec x).\phi)(\vec a)} &= \pi\cdot(\lfp(\lambda f.\lambda\vec b. \sem[{\xi[X\mapsto f]}]{\phi\{\vec b/\vec x\}}))(\vec a)\\
	&= (\lfp(\lambda f.\lambda\vec b. \pi\cdot\sem[{\xi[X\mapsto (\pi^{-1}\cdot f)]}]{\phi\{\pi^{-1} \cdot \vec b/\vec x\}}))(\pi\cdot \vec a)\\
	&= (\lfp(\lambda f.\lambda\vec b. \semp[\pi\cdot(\xi{[X\mapsto \pi^{-1}\cdot f]})]{\phi\{\pi^{-1}\cdot \vec b / \vec x\}}))(\pi\cdot \vec a)\ (IH)\\
	&= (\lfp(\lambda f.\lambda\vec b. \sem[(\pi\cdot\xi){[X\mapsto f]})]{(\pi \cdot\phi)\{\vec b / \vec x\}}))(\pi\cdot \vec a)\\
	&= \sem[\pi\cdot\xi]{(\mu X(\vec x).(\pi\cdot\phi)(\vec a)}\\
	&= \semp{(\mu X(\vec x).\phi(\vec a)}
\end{align*}
		
	\end{itemize}
\end{proof}
%The proof is similar to Lemma~\ref{lem:equivariant}, except for the following additional rule. 
%We show the new inductive step as follows:
%\begin{itemize}
%
%\end{itemize}
\subsection{Proof of \cref{thm:semanticswelldefined}}
\semanticswelldefined*
%	The semantics of \cref{def:hdsemantics} is well defined.
%	 Its image is within the subset of $\Ucal$ of finitely-supported sets of configurations.
%\end{theorem}
\begin{proof}
For the former, we must show that least fixpoints exist for any $\lfp(\lambda f.\lambda \vec b.\sem[{\xi[X \mapsto f]}]{\phi\{\vec b/\vec x\}})$. 
By definition, the function $F = \lambda f.\lambda \vec b.\sem[{\xi[X \mapsto f]}]{\phi\{\vec b/\vec x\}}$ is in $\Ucal_{\ar(X)}$ and by Lemma~\ref{lem:completelattice}, $\Ucal_{\ar(X)}$ is a complete lattice and $F$ is monotone. 
Thus, by the Knaster-Tarski theorem~\cite{Tarski}, the least fixpoint of any such $F$ exists.
\\
The proof of the latter follows from \cref{lem:semanticsequivariant}.
%\ntnote{only the second part follows from Lemma 13, not the well-defined part}
\end{proof}
%%% Local Variables:
%%% mode: latex
%%% TeX-master: "main"
%%% End:
